% Syllabus Units IV-V topics not on this year's slides: PGP, S/MIME, IPSec, SSL/TLS, HTTPS, SET, intruders, malware,
% firewalls, trusted systems. Sources: Stallings (CNS, chs. on e-mail, IP, web and system security), Forouzan chs. 16-18.
% Every number (radix-64, AH/ESP sizes, anti-replay trace, TLS record sizes, dual-signature toy) checked by scratchpad/w24/net.py.
\usetikzlibrary{decorations.pathreplacing,shapes.symbols}
\tikzset{
  nsB/.style={draw, fill=acc!8, font=\footnotesize, minimum height=0.55cm, inner sep=2pt, align=center},
  nsC/.style={draw, circle, inner sep=0pt, minimum size=4.6mm, font=\scriptsize},
  nsL/.style={font=\footnotesize, inner sep=1pt},
  nsA/.style={-{Stealth[length=1.6mm]}},
  nsBr/.style={decorate, decoration={brace, amplitude=3pt}}
}
\providecommand{\nsrect}[6]{\draw[fill=#6] (#1,#2) rectangle ++(#3,#4); \node[font=\scriptsize, align=center, inner sep=0pt] at (#1+#3/2,#2+#4/2) {#5};}% x, y, w, h, label, fill

\section{Network Security Protocols and System Security}\label{sec:netsec}

The cryptographic tools of the earlier sections (block ciphers and modes, RSA and Diffie--Hellman, hashes, MACs and signatures) become useful only when a protocol fixes exactly which tool protects which bytes, how the keys get there, and in what order the operations are applied. This part of the syllabus studies the standard protocols at three layers of the TCP/IP stack, e-mail security (PGP, S/MIME), IP security (IPSec), web security (SSL/TLS, HTTPS, SET), and then the defence of the computer systems themselves (intruders, malicious software, firewalls, trusted systems). None of it is on this year's slides, so the notes follow Stallings and Forouzan.

\subsection{Where security can be placed in the stack \new}
\likely{``Compare security at the network, transport and application layers'' (3 marks).}

Security can be added at three places, and the choice decides who must change and what is protected.
\begin{center}
\begin{tikzpicture}[font=\footnotesize]
\tikzset{lay/.style={draw, minimum width=2.9cm, minimum height=0.5cm, inner sep=1pt}}
\node[lay] at (0,1.0) {HTTP\ \ FTP\ \ SMTP}; \node[lay] at (0,0.5) {TCP}; \node[lay, fill=acc!15] at (0,0) {IP / IPSec};
\node[lay] at (4,1.5) {HTTP\ \ FTP\ \ SMTP}; \node[lay, fill=acc!15] at (4,1.0) {SSL or TLS}; \node[lay] at (4,0.5) {TCP}; \node[lay] at (4,0) {IP};
\node[lay, fill=acc!15] at (8,1.5) {S/MIME\ \ PGP\ \ SET}; \node[lay] at (8,1.0) {SMTP\ \ HTTP}; \node[lay] at (8,0.5) {TCP}; \node[lay] at (8,0) {IP};
\foreach \x/\d in {0/{(a) network level}, 4/{(b) transport level}, 8/{(c) application level}} \node[font=\footnotesize\itshape] at (\x,-0.6) {\d};
\end{tikzpicture}
\end{center}
\textbf{Network level (IPSec).} Security is applied to every IP packet, so it is transparent to applications and to users: no application has to be rewritten, and all traffic between two sites (web, mail, file transfer) is protected at once. \textbf{Transport level (SSL/TLS).} Security sits just above TCP and is used by the applications that ask for it (a browser opening an \ttx{https} connection); it protects one connection at a time. \textbf{Application level (PGP, S/MIME, SET).} Security is built into the application and can be tailored to it: an e-mail is signed by its author and stays signed in the recipient's mailbox, and SET can show the bank the card number while hiding it from the merchant. The price is that every application must implement its own security.

\begin{center}\small
\begin{tabularx}{\textwidth}{@{}l *{3}{>{\raggedright\arraybackslash}X}@{}}
\toprule
 & Network (IPSec) & Transport (SSL/TLS) & Application (PGP, S/MIME, SET)\\
\midrule
Protects & every IP packet between two hosts or gateways & one TCP connection & one message or transaction\\
Transparent to applications & yes & no (the application calls TLS) & no (built into the application)\\
Protection after delivery & none (removed on arrival) & none & yes (a signed mail stays signed)\\
Typical use & VPNs, site-to-site links & HTTPS, secure mail transport & signed/encrypted e-mail, card payments\\
\bottomrule
\end{tabularx}
\end{center}

\section{E-mail Security: PGP and S/MIME}\label{sec:email}

\subsection{Why e-mail needs its own security \new}
\likely{``Why can e-mail not use a handshake like TLS?'' (2 marks), as part of a PGP question.}

E-mail is a \textbf{one-way, store-and-forward} service: Alice writes the message, it is passed from server to server, and Bob reads it hours later. Since Alice and Bob are never online together, there is no handshake in which they could agree on algorithms or a session key. Alice must therefore decide everything herself and put it \emph{into} the message: the identifiers of the algorithms, a fresh session key encrypted for Bob, her signature and the information Bob needs to find her public key. Two practical constraints of the mail system shape the design as well: classic SMTP carries only 7-bit ASCII text, so binary ciphertext must be converted to printable characters, and mail servers limit the size of a message. PGP and S/MIME are the two standard solutions; both use the same cryptographic idea (hybrid encryption plus a hash-and-sign signature) and differ mainly in formats and in how public keys are trusted.

\subsection{PGP: the five services \hot\new}
\likely{``Explain the operational description of PGP with a diagram (authentication, confidentiality, both)'' (5--8 marks).}

\textbf{Pretty Good Privacy} (PGP), written by Phil Zimmermann in 1991, is a freely available package that provides confidentiality and authentication for e-mail and file storage. It selected the best available algorithms as building blocks, runs on every platform, and is not controlled by any government or standards body, which explains its spread. PGP offers five services.
\begin{center}\small
\begin{tabularx}{\textwidth}{@{}l l >{\raggedright\arraybackslash}X@{}}
\toprule
Service & Algorithms (classic PGP) & What is done\\
\midrule
Digital signature & DSS/SHA-1 or RSA/SHA-1 & a SHA-1 hash of the message is signed with the sender's private key and attached\\
Message encryption & CAST-128, IDEA or 3DES; session key by RSA or ElGamal & the message is encrypted with a one-time session key, which is encrypted with the recipient's public key and attached\\
Compression & ZIP & the message is compressed for storage and transmission\\
E-mail compatibility & radix-64 conversion & the binary result is converted to printable ASCII\\
Segmentation & --- & long messages are split to fit mail size limits and reassembled\\
\bottomrule
\end{tabularx}
\end{center}

\textbf{Authentication only.} The sender creates the message $M$ and computes its 160-bit SHA-1 hash $H(M)$. The hash is encrypted with RSA under the sender's private key $PR_a$ and prepended to the message. The receiver decrypts the signature with the sender's public key $PU_a$ to recover the hash, computes $H(M)$ of the received message himself, and accepts if the two are equal. Since only Alice holds $PR_a$, a match proves that she created the signature, and since the hash is collision resistant (\S\ref{sec:integrity}), it proves that $M$ was not changed. Signing the hash instead of $M$ keeps the RSA operation short whatever the length of the message.

\textbf{Confidentiality only.} The sender generates a random 128-bit \textbf{session key} $K_s$, used for this message only, and encrypts the message with CAST-128 (or IDEA or 3DES) in 64-bit CFB mode (\S\ref{sec:modes}). The session key is encrypted with RSA under the recipient's public key $PU_b$ and prepended. The receiver recovers $K_s$ with his private key $PR_b$ and decrypts the message. This is \textbf{hybrid encryption}: the slow public-key operation is applied to 128 bits only, the bulk is encrypted by the fast symmetric cipher, and the one-time key removes the need for any earlier key exchange, which a store-and-forward service could not have.

\textbf{Both services.} The sender first signs the plaintext, then encrypts the message together with the signature under $K_s$, and attaches $E(PU_b, K_s)$:
\[
A \to B:\quad E(PU_b, K_s) \,\big\|\, E_{K_s}\!\Big(Z\big(E(PR_a, H(M)) \,\|\, M\big)\Big),
\]
where $Z$ is ZIP compression; the whole result is then radix-64 encoded.

\begin{center}
\begin{tikzpicture}[font=\footnotesize]
% sender
\node[nsL, font=\footnotesize\bfseries, anchor=west] at (-0.6,1.85) {Source A};
\node[nsL] (m) at (-0.3,0) {$M$};
\node[nsB] (h) at (0.9,1.0) {$H$};
\node[nsB] (ep1) at (2.1,1.0) {$E_P$};
\node[nsL] (pra) at (2.1,1.75) {$PR_a$};
\node[nsC] (c1) at (3.2,0) {$\|$};
\node[nsB] (z) at (4.2,0) {$Z$};
\node[nsB] (ec) at (5.4,0) {$E_C$};
\node[nsL] (ks) at (5.4,1.0) {$K_s$ (random)};
\node[nsB] (ep2) at (6.9,1.0) {$E_P$};
\node[nsL] (pub) at (6.9,1.75) {$PU_b$};
\node[nsC] (c2) at (7.9,0) {$\|$};
\node[nsB] (r64) at (9.0,0) {R64};
\node[nsL] (net) at (10.4,0) {to mail};
\draw[nsA] (m) -- (c1); \draw[nsA] (0.1,0) |- (h); \draw[nsA] (h) -- (ep1); \draw[nsA] (pra) -- (ep1);
\draw[nsA] (ep1) -| (c1); \draw[nsA] (c1) -- (z); \draw[nsA] (z) -- (ec); \draw[nsA] (ks) -- (ec);
\draw[nsA] (ks) -- (ep2); \draw[nsA] (pub) -- (ep2); \draw[nsA] (ep2) -| (c2); \draw[nsA] (ec) -- (c2);
\draw[nsA] (c2) -- (r64); \draw[nsA] (r64) -- (net);
% receiver
\begin{scope}[yshift=-3.3cm]
\node[nsL, font=\footnotesize\bfseries, anchor=west] at (-0.6,1.85) {Destination B};
\node[nsL] (in) at (-0.3,0) {from mail};
\node[nsB] (r64i) at (1.0,0) {R64$^{-1}$};
\node[nsC] (s1) at (2.1,0) {$\div$};
\node[nsB] (dp1) at (3.1,1.0) {$D_P$};
\node[nsL] (prb) at (3.1,1.75) {$PR_b$};
\node[nsB] (dc) at (4.3,0) {$D_C$};
\node[nsB] (zi) at (5.4,0) {$Z^{-1}$};
\node[nsC] (s2) at (6.4,0) {$\div$};
\node[nsB] (dp2) at (7.7,1.0) {$D_P$};
\node[nsL] (pua) at (7.7,1.75) {$PU_a$};
\node[nsB] (h2) at (7.7,0) {$H$};
\node[nsC] (cmp) at (9.0,0.5) {$=$?};
\node[nsL] (mo) at (9.2,-0.9) {$M$};
\draw[nsA] (in) -- (r64i); \draw[nsA] (r64i) -- (s1); \draw[nsA] (s1) |- (dp1); \draw[nsA] (prb) -- (dp1);
\draw[nsA] (dp1) -| node[nsL, right, pos=0.7]{$K_s$} (dc); \draw[nsA] (s1) -- (dc); \draw[nsA] (dc) -- (zi); \draw[nsA] (zi) -- (s2);
\draw[nsA] (s2) |- (dp2); \draw[nsA] (pua) -- (dp2); \draw[nsA] (dp2) -| (cmp);
\draw[nsA] (s2) |- (6.8,-0.9) -- (mo); \draw[nsA] (6.95,-0.9) |- (h2); \draw[nsA] (h2) -| (cmp);
\end{scope}
\end{tikzpicture}
\end{center}
$E_P$/$D_P$ are public-key encryption/decryption (RSA), $E_C$/$D_C$ the symmetric cipher, $Z$ compression, R64 radix-64 conversion, $\|$ concatenation and $\div$ separation of the components.

At the receiver the steps run in reverse: radix-64 decoding, recovery of $K_s$ with $PR_b$, symmetric decryption, decompression, and finally verification of the signature with $PU_a$ against a freshly computed $H(M)$.

\subsection{PGP: the order of operations and why \hot\new}
\likely{``Why does PGP sign before compressing and compress before encrypting?'' (3 marks).}

The sender applies the operations in the order \textbf{sign $\to$ compress $\to$ encrypt $\to$ radix-64 $\to$ segment}, and each position has a reason.

\textbf{Sign before compressing.} First, the signature is then computed on the uncompressed message, so the recipient can store the plain message together with the signature and verify it later without keeping a compressed copy or recompressing. Second, PGP's compression is not deterministic: different versions of ZIP trade speed against compression ratio and produce different compressed outputs from the same input. A signature on the compressed form would verify only with the same compressor version; a signature on the original message is independent of the compressor.

\textbf{Compress before encrypting.} Compression removes redundancy, and redundancy is what cryptanalysis feeds on, so encrypting the compressed text strengthens the encryption. It also shortens what has to be encrypted and sent. (Compressing \emph{after} encryption would be useless anyway, since good ciphertext looks random and does not compress.)

\textbf{Radix-64 after encrypting.} The signature, the encrypted session key and the ciphertext are arbitrary binary data, which 7-bit mail systems may corrupt. Radix-64 (the same encoding as MIME base64) maps every 3 bytes to 4 printable characters from a 64-symbol alphabet (A--Z, a--z, 0--9, $+$, $/$), and a 24-bit CRC is appended to detect transmission errors. It expands the data by $4/3$, about 33\%, but ZIP compresses ordinary text by about a factor of 2, so the final message is usually \emph{shorter} than the plaintext: $\frac{1}{2} \times \frac{4}{3} \approx 0.67$.

\textbf{Segmentation last.} Many mail systems limit the message size (e.g.\ 50\,000 bytes), so after all other processing PGP splits a long message into segments that are mailed separately. The session key component and the signature component appear only once, at the start of the first segment; the receiver strips the mail headers and reassembles the whole block before any other processing.

\begin{worked}[: radix-64 and the size of a PGP message]
\textbf{Q.} (a) Radix-64 encode the three bytes \ttx{4D 61 6E} (``Man''). (b) A 60\,000-byte text is signed (signature component 300 bytes), compressed by ZIP with ratio 2, encrypted (CFB, no expansion), given a 260-byte session key component, and radix-64 encoded. Ignoring line breaks and the CRC, how many characters are sent, and how many segments if the mailer accepts at most 20\,000 characters?

\textbf{A.} (a) The 24 bits $01001101\ 01100001\ 01101110$ are regrouped into four 6-bit groups $010011\ 010110\ 000101\ 101110 = 19, 22, 5, 46$. In the radix-64 alphabet (0--25 $\to$ A--Z, 26--51 $\to$ a--z, 52--61 $\to$ 0--9) these are \textbf{T, W, F, u}, so the output is \ttx{TWFu}.\\
(b) Signed: $60\,000 + 300 = 60\,300$ bytes; compressed: $60\,300 / 2 = 30\,150$; with the session key component: $30\,410$ bytes; radix-64: $\lceil 30\,410 / 3 \rceil \times 4 = 10\,137 \times 4 = \mathbf{40\,548}$ \textbf{characters}; segments: $\lceil 40\,548 / 20\,000 \rceil = \mathbf{3}$.\\
\emph{Check}: $40\,548 / 60\,000 = 0.676$, close to the predicted $\frac12 \cdot \frac43 = 0.67$.
\end{worked}

\textbf{Message format.} A PGP message consists of three components. The \emph{message component} holds the data with a filename and a timestamp. The \emph{signature component} holds a timestamp, the \textbf{key ID} of the sender's public key, the \textbf{leading two octets of the message digest} in the clear, and the digest encrypted with the sender's private key. The two clear octets let the receiver check, before doing anything else, that he used the right public key to decrypt the digest: if the decrypted digest does not start with them, the wrong key was used. The \emph{session key component} holds the key ID of the recipient's public key and the encrypted session key. Compression and symmetric encryption cover the signature and message components; radix-64 covers everything.

\subsection{PGP: key IDs and key rings \new}
\likely{``Explain the PGP private-key ring and public-key ring; what is a key ID?'' (4 marks).}

\textbf{Key IDs.} A user may own several key pairs (to change keys from time to time, or to separate work and private use), and he will hold the public keys of many correspondents. The recipient therefore needs to know which public key was used. Sending the whole public key would waste space, so PGP sends a \textbf{key ID}, the least significant 64 bits of the public key: $\text{ID} = PU \bmod 2^{64}$. With 64 bits, the chance that two keys of one user share an ID is negligible.

\textbf{Private-key ring.} Each user stores his own key pairs in a table with one row per pair: timestamp (when the pair was created), key ID, public key, \textbf{encrypted private key}, and user ID (normally the e-mail address); the table can be indexed by user ID or key ID. The private key is never stored in the clear. When a pair is generated, PGP asks for a passphrase, computes its 160-bit SHA-1 hash, and encrypts the private key with CAST-128 using 128 bits of that hash as the key; the passphrase and the hash are then discarded. To sign, the user types the passphrase again, PGP recomputes the hash and decrypts the private key. A thief who copies the ring therefore still needs the passphrase, so the security of the private key rests on the passphrase.

\textbf{Public-key ring.} The public keys of other users are stored with timestamp, key ID, public key, user ID and the three trust fields described next: \textbf{owner trust}, the \textbf{signatures} on the key with their \textbf{signature trust}, and the computed \textbf{key legitimacy}. When a message arrives, PGP uses the key ID in the session key component to fetch the right private key from the private-key ring, and the key ID in the signature component to fetch the sender's public key from the public-key ring.

\subsection{PGP: the web of trust \new}
\likely{``Explain the PGP trust model: owner trust, signature trust, key legitimacy'' (4 marks), with a small example.}

\textbf{Problem.} Public-key encryption is only as safe as the binding between a public key and its owner. If Eve can plant a key with Bob's user ID on Alice's ring, Alice will encrypt to Eve and accept Eve's signatures as Bob's (a man-in-the-middle). PGP has no central certificate authority. Instead, any user may sign another user's public key to vouch that it belongs to that user (a key signature acts as a certificate), and each user decides whom he trusts as an \textbf{introducer}. The resulting graph of mutual signatures is the \textbf{web of trust}.

\textbf{Mechanism.} Each entry of Alice's public-key ring carries three kinds of trust information.
\begin{itemize}
  \item \textbf{Owner trust} says how far Alice trusts the owner of this key to sign \emph{other} people's keys correctly. It is assigned by Alice herself: unknown, usually not trusted, usually (partially) trusted, always (fully) trusted, or ultimate for her own keys.
  \item \textbf{Signature trust} is attached to each signature on the key and is simply a copy of the owner trust of the signer, if the signer is on Alice's ring (otherwise unknown).
  \item \textbf{Key legitimacy} says how far Alice believes the key really belongs to its user ID. PGP computes it from the signature trusts: one signature of ultimate trust makes it complete; otherwise PGP adds weights, $1/X$ for each always-trusted signature and $1/Y$ for each usually-trusted one ($X$ and $Y$ are user parameters, typically $X = 1$, $Y = 2$), and when the sum reaches 1 the key is complete, below 1 it is marginal.
\end{itemize}
The important distinction is that \textbf{legitimacy} concerns the key (is it really Bob's?) while \textbf{owner trust} concerns the person (does Bob check keys carefully before signing them?). A key can be fully legitimate while its owner is not trusted as an introducer. A user revokes a compromised key by issuing a \textbf{key revocation certificate} signed with the corresponding private key and spreading it as widely as possible.

\begin{worked}[: key legitimacy in a web of trust]
\textbf{Q.} Alice gives owner trust \emph{always} to Bob and \emph{usually} to Carol and Dave, and has not assigned any trust to Frank. Weights are 1 for an always-trusted and $\frac12$ for a usually-trusted signer. Determine the legitimacy of the keys of (a) Gita, signed by Bob; (b) Hari, signed by Carol and Dave; (c) Ivan, signed by Carol; (d) Jaya, signed by Frank, whose own key on Alice's ring is legitimate.

\textbf{A.} (a) $1 \ge 1$: \textbf{complete}. (b) $\frac12 + \frac12 = 1$: \textbf{complete}. (c) $\frac12 < 1$: \textbf{marginal}. (d) Frank's owner trust is unknown, so his signature has weight 0: Jaya's key is \textbf{not legitimate}, although Frank's own key is. Legitimacy of the signer's key is not enough; Alice must trust him as an introducer.\\
\emph{Check}: only (a) and (b) reach the threshold 1.
\end{worked}

\subsection{S/MIME and MIME \new}
\likely{``What are the limitations of SMTP/RFC 822 that MIME removes? List the MIME header fields and content types'' (4 marks).}

\textbf{Problem.} The original Internet mail format (RFC 822) and SMTP were designed for short 7-bit ASCII text. As a result, SMTP cannot carry executable files or other binary objects, cannot carry national-language characters (8-bit codes), and may reject messages above a size limit; gateways translating between ASCII and EBCDIC or to X.400 networks use inconsistent mappings; and non-standard SMTP implementations delete or add carriage returns and line feeds, truncate or wrap lines longer than 76 characters, remove trailing white space, pad lines to the same length and convert tabs to spaces. Any of these silently corrupts binary data such as ciphertext or a signature.

\textbf{MIME} (Multipurpose Internet Mail Extensions) solves this within RFC 822 by adding five header fields (MIME-Version, Content-Type, Content-Transfer-Encoding, Content-ID and Content-Description), a set of standard content types, and transfer encodings (7bit, 8bit, binary, quoted-printable for mostly-ASCII text, and \textbf{base64} for arbitrary binary data) that make any content survive the mail system.
\begin{center}\small
\begin{tabularx}{\textwidth}{@{}l l >{\raggedright\arraybackslash}X@{}}
\toprule
Type & Subtypes & Content\\
\midrule
text & plain, enriched & unformatted or formatted text\\
multipart & mixed, parallel, alternative, digest & several parts in one message (in order, in any order, alternative versions, digest of messages)\\
message & rfc822, partial, external-body & an encapsulated message, a fragment of a large one, a pointer to one\\
image, video, audio & jpeg, gif; mpeg; basic & media\\
application & postscript, octet-stream & other data, e.g.\ arbitrary binary\\
\bottomrule
\end{tabularx}
\end{center}

\textbf{S/MIME} (Secure/MIME), based on RSA Data Security's technology and standardised by the IETF, adds security to MIME by defining new content types that carry signed and/or encrypted data in the PKCS\#7 (Cryptographic Message Syntax) format. It provides four functions.
\begin{itemize}
  \item \textbf{Enveloped data} (\ttx{application/pkcs7-mime; smime-type=envelopedData}): the sender generates a random session key, encrypts the content with it (3DES or AES), and for \emph{each} recipient encrypts the session key with that recipient's public RSA key and builds a RecipientInfo block (recipient's certificate identifier, algorithm identifier, encrypted session key). The RecipientInfo blocks and the encrypted content are then base64 encoded. This gives confidentiality, and one message serves many recipients with one bulk encryption.
  \item \textbf{Signed data} (\ttx{smime-type=signedData}): the sender hashes the content (SHA), encrypts the digest with his private key, and builds a SignerInfo block (signer's certificate, digest and signature algorithm identifiers, encrypted digest). Content and signature are base64 encoded \emph{together}, so only an S/MIME-capable reader can see the content.
  \item \textbf{Clear-signed data} (\ttx{multipart/signed}): the signature is computed the same way, but the message has two parts: the content in its normal MIME form, readable by anyone, and a detached signature (\ttx{application/pkcs7-signature}) in base64. Recipients without S/MIME can still read the message, though they cannot verify it.
  \item \textbf{Signed and enveloped data}: the two operations can be nested in either order, signing the plaintext and then enveloping it, or enveloping and then signing the ciphertext.
\end{itemize}

\textbf{Certificates.} S/MIME uses \textbf{X.509 version 3} certificates issued by certification authorities. Its key management is a hybrid of a strict CA hierarchy and PGP's web of trust: each client is configured with a list of trusted root certificates and certificate revocation lists, and a user agent must generate key pairs, request certificates from a CA, and store and retrieve certificates of its correspondents.

\subsection{PGP versus S/MIME \new}
\likely{``Differentiate PGP and S/MIME'' (3--4 marks).}
\begin{center}\small
\begin{tabularx}{\textwidth}{@{}l *{2}{>{\raggedright\arraybackslash}X}@{}}
\toprule
 & PGP & S/MIME\\
\midrule
Origin & one author (Zimmermann), free software; OpenPGP standard & RSA Data Security / IETF standard, built into mail clients\\
Trust model & web of trust: users sign each other's keys & hierarchical CAs with X.509v3 certificates\\
Certificate format & PGP's own key certificates & X.509v3\\
Message format & PGP's own binary format & MIME with PKCS\#7 / CMS content types\\
Binary-to-text encoding & radix-64 with CRC & base64 (MIME transfer encoding)\\
Compression & ZIP, built in & optional (compressed-data type)\\
Readable without the software & no & yes for clear-signed messages\\
Typical users & individuals & organisations and industry\\
\bottomrule
\end{tabularx}
\end{center}
Both use the same cryptographic pattern: hash-and-sign with the sender's private key, then hybrid encryption with a one-time session key wrapped in the recipient's public key.

\section{IP Security (IPSec)}\label{sec:ipsec}

\subsection{Why security at the IP level; applications and services \new}
\likely{``What are the applications and benefits of IPSec? List the services provided by AH and ESP'' (4 marks).}

\textbf{Problem.} IP itself offers no security: the source address of a packet can be forged, and any router or eavesdropper on the path can read or change the contents. Securing each application separately (as PGP does for mail) leaves all other traffic unprotected. \textbf{IPSec} adds authentication, integrity and confidentiality to IP itself, so that every packet between two hosts or two networks can be protected whatever application produced it. It is optional in IPv4 and was designed together with IPv6.

\textbf{Applications.} IPSec is used for secure branch-office connectivity over the Internet (a \textbf{virtual private network}, VPN, instead of leased lines), secure remote access for employees through an ISP, extranet and intranet connections with partners, and as a basis for electronic commerce, since all traffic between the sites is protected without changes to the applications.

\textbf{Benefits.} When IPSec is implemented in a firewall or router, it protects all traffic crossing the perimeter, while traffic inside the site carries no security overhead; and if all outside traffic must use IP through the firewall, IPSec there cannot be bypassed. Because it lies below the transport layer, it is \textbf{transparent to applications and to end users}: no software on the hosts has to change and users need no training or keys. It can also protect individual users where needed, and it secures routing (router advertisements and redirects really come from authorised routers).

\textbf{Architecture and services.} IPSec consists of two traffic protocols, the \textbf{Authentication Header} (AH), which authenticates packets, and the \textbf{Encapsulating Security Payload} (ESP), which encrypts and optionally authenticates them, plus a key-management protocol (IKE) and the documents that define the algorithms. Together they provide the following services.
\begin{center}\small
\begin{tabular}{@{}l c c c@{}}
\toprule
Service & AH & ESP (encryption only) & ESP (encryption $+$ authentication)\\
\midrule
Access control & \checkmark & \checkmark & \checkmark\\
Connectionless integrity & \checkmark & & \checkmark\\
Data origin authentication & \checkmark & & \checkmark\\
Rejection of replayed packets & \checkmark & \checkmark & \checkmark\\
Confidentiality & & \checkmark & \checkmark\\
Limited traffic-flow confidentiality & & \checkmark & \checkmark\\
\bottomrule
\end{tabular}
\end{center}

\subsection{Security associations, SAD and SPD \hot\new}
\likely{``What is a security association? What parameters identify it? Explain SAD and SPD'' (4 marks).}

\textbf{Idea.} Before two parties can protect packets they must agree on the protocol (AH or ESP), the algorithms, the keys, the lifetime and so on. IPSec stores such an agreement as a \textbf{security association} (SA): \imp{a one-way relationship between a sender and a receiver that affords security services to the traffic carried on it.} Since it is one-way, a two-way conversation needs two SAs, one in each direction, and an SA covers one protocol only, AH or ESP. An SA is uniquely identified by three parameters:
\begin{itemize}
  \item the \textbf{Security Parameters Index} (SPI), a 32-bit value chosen by the receiver and carried in every AH and ESP header, which lets the receiver select the SA for an arriving packet;
  \item the \textbf{IP destination address}, the address of the SA's endpoint (an end host or a firewall/router);
  \item the \textbf{security protocol identifier}, which says whether the SA is AH or ESP.
\end{itemize}

\textbf{Security Association Database (SAD).} Each IPSec host keeps one entry per SA with the parameters needed to process its packets: the sequence number counter and its overflow flag, the anti-replay window, the AH information (authentication algorithm and keys), the ESP information (encryption and authentication algorithms, keys, IVs), the lifetime of the SA, the IPSec mode (transport or tunnel) and the path MTU.

\textbf{Security Policy Database (SPD).} The SPD decides \emph{which} traffic gets \emph{which} protection. Each entry is defined by \textbf{selectors}, values of IP and upper-layer fields (destination and source IP addresses, user ID, data sensitivity level, transport protocol, source and destination ports), and specifies an action: \textbf{discard} the packet, \textbf{bypass} IPSec, or \textbf{protect} it with the SA(s) the entry points to. For every outbound packet IPSec compares the packet's fields with the SPD, finds the matching entry, determines the SA (creating one with IKE if none exists yet) and applies its processing. For an inbound packet the receiver uses the SPI, destination address and protocol to find the SA in the SAD, processes the packet, and then checks that the result is allowed by the SPD.

\subsection{Transport mode and tunnel mode \hot\new}
\likely{``Differentiate transport and tunnel mode; draw the packet formats for AH and ESP in both modes'' (5 marks).}

Both AH and ESP can be used in two modes, which differ in what is protected and in where the IPSec processing is done.

\textbf{Transport mode} protects the \emph{payload} of the IP packet, i.e.\ the upper-layer segment (TCP or UDP header plus data). The AH or ESP header is inserted between the original IP header and the payload, and the original IP header stays visible. It is used for \textbf{end-to-end} communication between two hosts that both run IPSec, such as a client and a server.

\textbf{Tunnel mode} protects the \emph{entire} original IP packet. After the AH or ESP fields are added, the whole packet (including its original header) becomes the payload of a \textbf{new outer IP packet} with a new header, normally addressed from one security gateway to another. Routers on the path see only the outer header, so the inner source and destination addresses are hidden. Tunnel mode is used when one or both ends of the SA are \textbf{security gateways} (firewalls or routers): hosts behind the gateways need no IPSec at all, and the gateways build a VPN between the sites.

\begin{center}
\begin{tikzpicture}[font=\scriptsize]
\def\h{0.5}
% (a) original
\node[anchor=east, align=right] at (-0.2,0.25) {original\\IPv4 packet};
\nsrect{0}{0}{1.6}{\h}{orig IP hdr}{gray!15}\nsrect{1.6}{0}{1.0}{\h}{TCP}{white}\nsrect{2.6}{0}{2.0}{\h}{Data}{white}
% (b) AH transport
\begin{scope}[yshift=-1.4cm]
\node[anchor=east, align=right] at (-0.2,0.25) {AH,\\transport};
\nsrect{0}{0}{1.6}{\h}{orig IP hdr}{gray!15}\nsrect{1.6}{0}{1.0}{\h}{AH}{gold!30}\nsrect{2.6}{0}{1.0}{\h}{TCP}{white}\nsrect{3.6}{0}{2.0}{\h}{Data}{white}
\draw[nsBr] (0,0.58) -- node[above=2pt]{authenticated except mutable fields} (5.6,0.58);
\end{scope}
% (c) AH tunnel
\begin{scope}[yshift=-2.8cm]
\node[anchor=east, align=right] at (-0.2,0.25) {AH,\\tunnel};
\nsrect{0}{0}{1.6}{\h}{new IP hdr}{acc!15}\nsrect{1.6}{0}{1.0}{\h}{AH}{gold!30}\nsrect{2.6}{0}{1.6}{\h}{orig IP hdr}{gray!15}\nsrect{4.2}{0}{1.0}{\h}{TCP}{white}\nsrect{5.2}{0}{2.0}{\h}{Data}{white}
\draw[nsBr] (0,0.58) -- node[above=2pt]{authenticated except mutable fields in the new IP header} (7.2,0.58);
\end{scope}
% (d) ESP transport
\begin{scope}[yshift=-4.6cm]
\node[anchor=east, align=right] at (-0.2,0.25) {ESP,\\transport};
\nsrect{0}{0}{1.6}{\h}{orig IP hdr}{gray!15}\nsrect{1.6}{0}{1.2}{\h}{ESP hdr}{ok!25}\nsrect{2.8}{0}{1.0}{\h}{TCP}{white}\nsrect{3.8}{0}{2.0}{\h}{Data}{white}\nsrect{5.8}{0}{1.2}{\h}{ESP trlr}{ok!25}\nsrect{7.0}{0}{1.2}{\h}{ESP auth}{ok!25}
\draw[nsBr] (1.6,0.58) -- node[above=2pt]{authenticated} (7.0,0.58);
\draw[nsBr] (7.0,-0.08) -- node[below=2pt]{encrypted} (2.8,-0.08);
\end{scope}
% (e) ESP tunnel
\begin{scope}[yshift=-6.4cm]
\node[anchor=east, align=right] at (-0.2,0.25) {ESP,\\tunnel};
\nsrect{0}{0}{1.6}{\h}{new IP hdr}{acc!15}\nsrect{1.6}{0}{1.2}{\h}{ESP hdr}{ok!25}\nsrect{2.8}{0}{1.6}{\h}{orig IP hdr}{gray!15}\nsrect{4.4}{0}{1.0}{\h}{TCP}{white}\nsrect{5.4}{0}{2.0}{\h}{Data}{white}\nsrect{7.4}{0}{1.2}{\h}{ESP trlr}{ok!25}\nsrect{8.6}{0}{1.2}{\h}{ESP auth}{ok!25}
\draw[nsBr] (1.6,0.58) -- node[above=2pt]{authenticated} (8.6,0.58);
\draw[nsBr] (8.6,-0.08) -- node[below=2pt]{encrypted} (2.8,-0.08);
\end{scope}
\end{tikzpicture}
\end{center}

\begin{center}\small
\begin{tabularx}{\textwidth}{@{}l *{2}{>{\raggedright\arraybackslash}X}@{}}
\toprule
 & Transport mode SA & Tunnel mode SA\\
\midrule
AH & authenticates the IP payload and selected (immutable) parts of the IP header & authenticates the entire inner IP packet and selected parts of the outer IP header\\
ESP & encrypts the IP payload & encrypts the entire inner IP packet (inner header included)\\
ESP with authentication & encrypts and authenticates the IP payload, but not the IP header & encrypts and authenticates the entire inner IP packet\\
Endpoints & hosts (end-to-end) & usually security gateways (VPN)\\
Addresses seen on the path & real source and destination & only the gateways' addresses\\
\bottomrule
\end{tabularx}
\end{center}

\subsection{Authentication Header (AH) \hot\new}
\likely{``Draw the AH format and explain each field; which IP header fields are authenticated?'' (4--5 marks).}

\textbf{Purpose.} AH provides data integrity and data origin authentication of IP packets, plus anti-replay protection. It prevents address spoofing, since the source address is authenticated, and it stops replay through the sequence number. It is based on a MAC (HMAC, \S\ref{sec:integrity}), so the two parties must share a secret key. AH provides \emph{no} confidentiality.

\begin{center}
\begin{tikzpicture}[font=\scriptsize]
\def\u{0.2}
\node[font=\footnotesize\bfseries] at (3.2,2.95) {AH};
\foreach \b/\x in {0/0, 8/1.6, 16/3.2, 31/6.2} \node at (\x+0.1,2.65) {\b};
\nsrect{0}{2.0}{1.6}{0.5}{Next header}{gold!20}\nsrect{1.6}{2.0}{1.6}{0.5}{Payload len}{gold!20}\nsrect{3.2}{2.0}{3.2}{0.5}{Reserved}{gold!20}
\nsrect{0}{1.5}{6.4}{0.5}{Security Parameters Index (SPI)}{white}
\nsrect{0}{1.0}{6.4}{0.5}{Sequence number}{white}
\nsrect{0}{0}{6.4}{1.0}{Authentication data (ICV), variable}{acc!10}
\begin{scope}[xshift=7.3cm]
\node[font=\footnotesize\bfseries] at (3.2,2.95) {ESP};
\foreach \b/\x in {0/0, 16/3.2, 24/4.8, 31/6.2} \node at (\x+0.1,2.65) {\b};
\nsrect{0}{2.0}{6.4}{0.5}{Security Parameters Index (SPI)}{white}
\nsrect{0}{1.5}{6.4}{0.5}{Sequence number}{white}
\nsrect{0}{0.5}{6.4}{1.0}{Payload data (variable)}{ok!15}
\nsrect{0}{0}{6.4}{0.5}{Padding (0--255 bytes)}{ok!15}
\nsrect{0}{-0.5}{3.2}{0.5}{padding (cont.)}{ok!15}\nsrect{3.2}{-0.5}{1.6}{0.5}{Pad length}{ok!15}\nsrect{4.8}{-0.5}{1.6}{0.5}{Next header}{ok!15}
\nsrect{0}{-1.3}{6.4}{0.8}{Authentication data (ICV), variable}{acc!10}
\draw[nsBr] (6.55,1.5) -- node[right=3pt]{encrypted} (6.55,-0.5);
\draw[nsBr] (7.85,2.5) -- node[right=3pt]{authenticated} (7.85,-0.5);
\end{scope}
\end{tikzpicture}
\end{center}

\textbf{Fields.}
\begin{itemize}
  \item \textbf{Next header} (8 bits) identifies the header that follows AH (e.g.\ TCP in transport mode, IP in tunnel mode).
  \item \textbf{Payload length} (8 bits) is the length of AH in 32-bit words, minus 2.
  \item \textbf{Reserved} (16 bits) for future use.
  \item \textbf{SPI} (32 bits) identifies the SA.
  \item \textbf{Sequence number} (32 bits) is a monotonically increasing counter used for anti-replay.
  \item \textbf{Authentication data} (variable, a multiple of 32 bits) holds the \textbf{integrity check value} (ICV), a MAC of the packet; the defaults were HMAC-MD5-96 and HMAC-SHA-1-96, i.e.\ the HMAC truncated to its first 96 bits.
\end{itemize}

\textbf{What the ICV covers.} The MAC must be the same at the sender and at the receiver, but some IP header fields legitimately change on the way. AH therefore computes the MAC over the IP header fields that are \textbf{immutable} in transit (version, header length, total length, identification, protocol, source address, destination address) or \textbf{mutable but predictable} at the endpoint (the destination address when source routing is used, which is set to its final value). The \textbf{mutable, unpredictable} fields (type of service, flags, fragment offset, time to live, header checksum) are set to zero for the calculation at both ends. The MAC also covers the AH header itself, with the authentication data field set to zero, and the entire upper-layer data, which is immutable. Because the source and destination addresses are inside the MAC, an attacker cannot change them without detection; this is also why AH does not work through NAT, which rewrites addresses.

\begin{worked}[: AH length field and packet size]
\textbf{Q.} AH uses HMAC-SHA-1-96. Find the value of the Payload length field and the size of the packet when a 37-byte TCP segment (20-byte header, 17 data bytes) is sent from a host with a 20-byte IPv4 header in transport mode and in tunnel mode.

\textbf{A.} Fixed AH fields: $1 + 1 + 2 + 4 + 4 = 12$ bytes; ICV $= 96$ bits $= 12$ bytes; AH $= 24$ bytes $= 6$ words, so Payload length $= 6 - 2 = \mathbf{4}$.\\
Transport: $20 + 24 + 37 = \mathbf{81}$ bytes. Tunnel: new IP header $+$ AH $+$ inner packet $= 20 + 24 + (20 + 37) = \mathbf{101}$ bytes.\\
\emph{Check}: tunnel mode costs exactly one extra 20-byte IP header.
\end{worked}

\subsection{Anti-replay service \new}
\likely{``Explain the IPSec anti-replay window with an example'' (3--4 marks).}

\textbf{Problem.} In a replay attack the attacker records an authenticated packet and sends it again later; the MAC is still valid, so authentication alone cannot reject it. IP gives no ordering guarantee, so the receiver also cannot simply demand increasing numbers: packets may legitimately arrive out of order.

\textbf{Sender side.} When an SA is created, the sender's sequence counter is set to 0; for each packet it is incremented and written into the Sequence number field, so the first packet carries 1. The counter must not wrap around past $2^{32} - 1$: before that the SA is terminated and a new SA with a new key is negotiated, so no number repeats under one key.

\textbf{Receiver side.} The receiver keeps a \textbf{window} of size $W$ (default 64) whose right edge $N$ is the highest sequence number received so far in a valid packet; for each number from $N - W + 1$ to $N$ it remembers whether a valid packet with that number has arrived (marked slot). For an arriving packet:
\begin{enumerate}
  \item if its number is \textbf{inside the window} and the slot is unmarked, the MAC is checked; if it is valid, the slot is marked;
  \item if its number is \textbf{to the right of the window}, the MAC is checked; if it is valid, the window is advanced so that this number becomes the right edge, and its slot is marked;
  \item if its number is \textbf{to the left of the window}, or the slot is already marked, or the MAC fails, the packet is discarded (and the event is logged).
\end{enumerate}
The MAC is checked before the window moves, so a forged packet with a huge sequence number cannot push the window forward.

\begin{worked}[: tracing the anti-replay window]
\textbf{Q.} $W = 8$ and the right edge is $N = 20$, so the window is 13--20; valid packets 13, 14, 16, 17, 19, 20 have arrived. Packets arrive with numbers 18, 16, 25, 15, 21 (all with valid MACs). What happens to each?
\begin{center}\small
\begin{tikzpicture}[font=\scriptsize]
\foreach \n [count=\i] in {13,...,20} {
  \pgfmathsetmacro\f{(\n==15 || \n==18) ? 0 : 1}
  \ifdim\f pt>0.5pt \nsrect{\i*0.6}{0}{0.6}{0.45}{\n}{acc!25}\else\nsrect{\i*0.6}{0}{0.6}{0.45}{\n}{white}\fi
}
\node[anchor=west] at (5.6,0.22) {window 13--20 (shaded $=$ received)};
\draw[nsA] (5.4,-0.25) -- node[below]{advance} (6.2,-0.25);
\end{tikzpicture}
\begin{tabular}{@{}c l l@{}}
\toprule
Packet & Position & Action and new window\\
\midrule
18 & inside, slot unmarked & accept, mark 18; window 13--20\\
16 & inside, slot marked & \textbf{replay}: discard\\
25 & right of window & accept; window slides to 18--25 (marked 18, 19, 20, 25)\\
15 & left of the new window & discard, although 15 was never received\\
21 & inside, slot unmarked & accept, mark 21\\
\bottomrule
\end{tabular}
\end{center}
\emph{Check}: packet 15 shows the cost of a finite window: a genuine but very late packet is dropped; a larger $W$ tolerates more reordering.
\end{worked}

\subsection{Encapsulating Security Payload (ESP) \hot\new}
\likely{``Draw the ESP format; why is padding needed? Compare AH and ESP'' (5 marks).}

\textbf{Purpose.} ESP provides \textbf{confidentiality} (encryption of the payload), limited traffic-flow confidentiality, and optionally the same authentication, integrity and anti-replay services as AH, all with one header and trailer (format above).

\textbf{Fields.} \textbf{SPI} and \textbf{sequence number} are as in AH. \textbf{Payload data} (variable) is the transport-layer segment (transport mode) or the whole inner IP packet (tunnel mode), encrypted; if the cipher needs an IV (e.g.\ AES-CBC, \S\ref{sec:modes}), it is carried in the clear at the start of this field. \textbf{Padding} (0--255 bytes) is followed by the \textbf{pad length} (8 bits, the number of padding bytes) and the \textbf{next header} (8 bits, the type of the encrypted payload). \textbf{Authentication data} (variable) is the ICV, computed over the ESP packet from the SPI to the next-header field, i.e.\ \emph{after} encryption.

\textbf{Why padding.} Padding is added for three reasons: a block cipher needs the plaintext (payload, padding, pad length and next header) to be a whole number of blocks; ESP requires the pad-length and next-header fields to end on a 32-bit boundary, so the ciphertext is a multiple of 4 bytes; and extra padding can hide the true length of the payload, which gives partial traffic-flow confidentiality.

\textbf{Encrypt-then-MAC.} Because the ICV is computed over the ciphertext, the receiver checks it first and discards bogus or replayed packets \emph{before} spending time on decryption, which reduces the impact of denial-of-service attacks, and the two operations can be done in parallel. Unlike AH, ESP's authentication does \emph{not} cover the IP header in front of the ESP header (in tunnel mode the inner header is covered, being inside the payload).

\begin{worked}[: ESP padding and packet size]
\textbf{Q.} ESP uses AES-CBC (16-byte blocks, 16-byte explicit IV) and HMAC-SHA-1-96 (12-byte ICV). The same 37-byte TCP segment as above is sent. Find the padding and the total packet size in transport and in tunnel mode (IPv4 header 20 bytes).

\textbf{A.} The encrypted part is payload $+$ padding $+$ pad length (1) $+$ next header (1), a multiple of 16.\\
Transport: payload $= 37$; $37 + 2 = 39$; padding $= (-39) \bmod 16 = \mathbf{9}$; encrypted part $= 48$ bytes. Total $= 20$ (IP) $+ 8$ (SPI, seq.) $+ 16$ (IV) $+ 48 + 12$ (ICV) $= \mathbf{104}$ bytes.\\
Tunnel: payload $=$ inner packet $= 20 + 37 = 57$; $57 + 2 = 59$; padding $= (-59) \bmod 16 = \mathbf{5}$; encrypted part $= 64$. Total $= 20 + 8 + 16 + 64 + 12 = \mathbf{120}$ bytes.\\
\emph{Check}: $48 = 3 \times 16$ and $64 = 4 \times 16$.
\end{worked}

\subsection{AH versus ESP \new}
\likely{``Differentiate AH and ESP'' (3 marks).}
\begin{center}\small
\begin{tabularx}{\textwidth}{@{}l *{2}{>{\raggedright\arraybackslash}X}@{}}
\toprule
 & AH & ESP\\
\midrule
Services & integrity, origin authentication, anti-replay & confidentiality, optional integrity/authentication, anti-replay, limited traffic-flow confidentiality\\
Encryption & none & payload (and inner header in tunnel mode)\\
Authentication covers & payload, AH header and immutable/predictable fields of the (outer) IP header & ESP header, payload and trailer only, not the IP header in front\\
Header/trailer & header only & header, trailer and ICV\\
Works through NAT & no (addresses are authenticated) & yes (transport/UDP encapsulation)\\
IP protocol number & 51 & 50\\
\bottomrule
\end{tabularx}
\end{center}
Since ESP with authentication offers almost everything AH does, plus confidentiality, ESP is by far the more widely used; AH is chosen only when the outer IP header itself must be authenticated or when encryption is not allowed.

\subsection{Combining security associations \new}
\likely{``Explain transport adjacency and iterated tunnelling; give the basic combinations of SAs'' (4 marks).}

One SA implements either AH or ESP, but not both, and a packet may need different services between different pairs of nodes (end-to-end authentication between hosts, plus encryption between the two gateways). IPSec therefore allows a packet to be processed by a sequence of SAs, called an \textbf{SA bundle}. SAs are combined in two ways. In \textbf{transport adjacency} more than one security protocol is applied to the same IP packet without tunnelling, for example ESP and then AH in transport mode between the same two hosts; since both act at the same endpoints, only one level of combination is possible. In \textbf{iterated tunnelling} several layers of protocols are applied by IP tunnelling, one inside the other; each tunnel may start and end at a different IPSec node along the path, so many levels of nesting are possible.

\textbf{Authentication plus confidentiality} can be obtained in three ways: ESP with its authentication option; transport adjacency with an inner ESP SA (no authentication) and an outer AH SA, which has the advantage that the authentication now also covers the IP addresses; or a bundle of an inner AH transport SA and an outer ESP tunnel SA, which authenticates the plaintext before encryption, so the authentication data are protected by encryption and can be stored with the message for later checking.

\textbf{Basic combinations} (Stallings). Case 1: end-to-end security between two IPSec hosts (AH or ESP in transport mode, or one inside the other). Case 2: security only between two gateways, a simple VPN with one tunnel SA, the hosts doing nothing. Case 3: case 2 plus case 1, a gateway-to-gateway tunnel that carries end-to-end transport SAs between the hosts. Case 4: a remote host reaches the organisation's firewall over the Internet with a tunnel SA, and an inner SA protects the traffic all the way to the server inside.

\subsection{Key management: Oakley, ISAKMP and IKE \hot\new}
\likely{``How does Oakley improve Diffie--Hellman? Explain cookies (clogging attack) and nonces (replay)'' (5 marks).}

\textbf{Problem.} AH and ESP need shared secret keys: two pairs per two-way SA bundle (one for integrity, one for encryption in each direction). With \textbf{manual} key management an administrator configures each system with its own and its peers' keys, which only works for small static networks. \textbf{Automated} management creates keys on demand; the default for IPSec is \textbf{ISAKMP/Oakley}, later called the Internet Key Exchange (\textbf{IKE}). Oakley is a key-determination protocol based on Diffie--Hellman (\S\ref{sec:keymgmt}); ISAKMP is the framework that defines the message formats and procedures for negotiating, establishing, modifying and deleting SAs.

\textbf{Weaknesses of plain Diffie--Hellman.} Diffie--Hellman lets two parties create a shared secret $g^{xy} \bmod p$ over a public channel, but (1) it provides no information about the identities of the parties, so a man in the middle can run one exchange with each side (\S\ref{sec:keymgmt}); and (2) it is computationally intensive, which opens a \textbf{clogging attack}: the attacker sends a flood of requests with forged source addresses, each forcing the victim to compute a modular exponentiation, until its processor is exhausted.

\textbf{Oakley's improvements.} Oakley keeps the advantages of Diffie--Hellman (keys are created only when needed, no pre-existing infrastructure beyond global parameters) and adds five features:
\begin{itemize}
  \item \textbf{Cookies} to thwart clogging attacks.
  \item Negotiation of a \textbf{group}, i.e.\ the global parameters of the exchange (768- or 1024-bit modular exponentiation, or elliptic-curve groups).
  \item \textbf{Nonces} to ensure against replay attacks.
  \item Exchange of the Diffie--Hellman public values.
  \item \textbf{Authentication} of the exchange, by digital signatures, public-key encryption or a symmetric key, to thwart man-in-the-middle attacks.
\end{itemize}

\textbf{Cookies.} Each side sends a pseudorandom number, the cookie, in its first message, and the other side must echo it in the first message of the Diffie--Hellman exchange. A cookie is sent to the claimed source address, so an attacker who forges the address never receives it and cannot echo it. As a result \imp{the attacker can only make the victim generate cookies, never perform the expensive exponentiation.} A cookie must satisfy three requirements: it must depend on the specific parties (so that a cookie obtained with the attacker's real address cannot be reused from forged addresses and ports); it must be impossible for anyone but the issuer to generate a cookie the issuer will accept (so it depends on a local secret, which cannot be deduced from any cookie); and generating and checking it must be fast. The recommended construction is a fast hash over the source and destination IP addresses, the source and destination ports and a locally generated secret value, $\mathrm{CKY} = H(\mathrm{IP}_s \| \mathrm{IP}_d \| \mathrm{port}_s \| \mathrm{port}_d \| \text{secret})$; the issuer stores nothing and simply recomputes the hash when the cookie comes back.

\begin{center}
\begin{tikzpicture}[font=\scriptsize]
\node[draw, fill=gray!10, font=\footnotesize] (i) at (0,0) {Initiator I};
\node[draw, fill=gray!10, font=\footnotesize] (r) at (12,0) {Responder R};
\draw[gray] (i) -- (0,-4.6); \draw[gray] (r) -- (12,-4.6);
\draw[nsA, thick] (0,-1.2) -- node[above, align=center]{\textbf{1}\; CKY$_I$, OK\_KEYX, GRP, $g^x$, EHAO, NIDP, ID$_I$, ID$_R$, $N_I$,\\$S_{KI}[\mathrm{ID}_I \| \mathrm{ID}_R \| N_I \| \mathrm{GRP} \| g^x \| \mathrm{EHAO}]$} (12,-1.2);
\draw[nsA, thick] (12,-2.6) -- node[above, align=center]{\textbf{2}\; CKY$_R$, CKY$_I$, OK\_KEYX, GRP, $g^y$, EHAS, NIDP, ID$_R$, ID$_I$, $N_R$, $N_I$,\\$S_{KR}[\mathrm{ID}_R \| \mathrm{ID}_I \| N_R \| N_I \| \mathrm{GRP} \| g^y \| g^x \| \mathrm{EHAS}]$} (0,-2.6);
\draw[nsA, thick] (0,-4.0) -- node[above, align=center]{\textbf{3}\; CKY$_I$, CKY$_R$, OK\_KEYX, GRP, $g^x$, EHAS, NIDP, ID$_I$, ID$_R$, $N_I$, $N_R$,\\$S_{KI}[\mathrm{ID}_I \| \mathrm{ID}_R \| N_I \| N_R \| \mathrm{GRP} \| g^x \| g^y \| \mathrm{EHAS}]$} (12,-4.0);
\end{tikzpicture}
\end{center}
\textbf{The aggressive Oakley exchange} (Stallings) above uses three messages. CKY are cookies, OK\_KEYX marks a key-exchange message, GRP names the Diffie--Hellman group, $g^x$ and $g^y$ are the public values, EHAO and EHAS are the encryption, hash and authentication algorithms offered and selected, NIDP says that the rest of the message is not encrypted, ID are identities, $N$ are nonces, and $S_K[\,\cdot\,]$ is a signature with the sender's private key.
\begin{itemize}
  \item \textbf{Message 1.} I sends its cookie, the group and its public value $g^x$, the algorithms it offers, both identities and a fresh nonce $N_I$, and signs the identities, nonce, group, $g^x$ and offer with its private key. The signature binds $g^x$ to I's identity, so a man in the middle cannot substitute his own value.
  \item \textbf{Message 2.} R returns its cookie and echoes I's, sends $g^y$, the selected algorithms and its nonce $N_R$, echoes $N_I$, and signs all of this \emph{including} $N_I$ and $g^x$. Signing I's fresh nonce proves that the reply is new and not a replay, and signing $g^x$ confirms which exchange it answers.
  \item \textbf{Message 3.} I echoes R's cookie, which proves that I really receives messages at its claimed address, and signs both nonces and both public values, which proves to R that I is live and holds the private key.
\end{itemize}
Both sides then compute $K = g^{xy} \bmod p$ and derive the keys from it. The aggressive mode saves messages but makes R compute $g^y$ before it has seen its own cookie echoed; the six-message \textbf{main mode} first exchanges cookies and the SA proposal (messages 1--2), then the Diffie--Hellman values and nonces (3--4), and only then the identities and authentication, encrypted under the new key (5--6), which also hides the identities from eavesdroppers.

\textbf{ISAKMP and the IKE phases.} An ISAKMP message starts with a header holding the initiator and responder cookies (64 bits each), the type of the next payload, version numbers, the exchange type, flags, a message ID and the length, followed by a chain of payloads (SA, proposal, transform, key exchange, identification, certificate, certificate request, hash, signature, nonce, notification, delete). IKE works in two phases. \textbf{Phase 1} (main or aggressive mode) authenticates the peers and builds a protected \textbf{IKE SA} between them. \textbf{Phase 2} (``quick mode'') runs under the protection of the IKE SA and quickly negotiates the actual AH/ESP SAs, and can be repeated to create more SAs or to rekey without a new Diffie--Hellman exchange. The current version, IKEv2, does the same with four messages (IKE\_SA\_INIT and IKE\_AUTH) and lets a busy responder answer the first message statelessly with a cookie, the same anti-clogging idea.
